\section{The last nonzero suspension and the sphere}\label{sec:conclusion}

Let \(\alpha\in\pi_d^S\), \(d>0\), have nonzero mod-2 Hurewicz
image. Represent it by \(f_0:S^d\to Q_0S^0\), and write
\[
f_j:S^{d+j}\longrightarrow QS^j,\qquad
z_j=h(f_j)=\sigma_*^j h(f_0)
\]
for the successive adjoints and their Hurewicz images. All the \(z_j\)
are primitive and \(A\)-annihilated. There is a largest \(e\ge0\)
with \(z_e\ne0\). Indeed, when \(j>d\), the degree \(d+j\) lies
strictly between \(j\) and \(2j\); the free infinite-loop homology
description has no class in that interval.

\begin{lemma}\label{lem:last-root}
For some \(m\ge2\),
\begin{equation}\label{eq:last-square}
z_e=w^2,\qquad |w|=m-1,\qquad d+e=2m-2,
\end{equation}
where \(w\) is primitive and \(A\)-annihilated, and
\(\sigma_*w\ne0\). If \(e\ge1\), then \(|w|\) is odd.
\end{lemma}

\begin{proof}
Since \(\sigma_*z_e=0\), Lemmas~\ref{lem:suspension-kernel}
and~\ref{lem:primitive-roots} give \eqref{eq:last-square} with a
primitive \(A\)-annihilated root. Its degree is positive, so \(m\ge2\).

If \(e=0\) and \(\sigma_*w=0\), the same lemmas make \(w\) a
square. Then \(z_0\) is a fourth power, contradicting
Proposition~\ref{prop:no-fourth-powers}. Thus the root suspends
nontrivially in this case.

For \(e\ge1\), we instead use the preceding spherical class
\(z_{e-1}\) to prove that \(|w|\) is odd. First suppose that
\(e\ge2\), so both levels have the direct weight decomposition
of Lemma~\ref{lem:weight-modules}. Choose a nonzero weight \(2^l\)
component of \(z_e\). Here \(l\ge1\), since the degree \(d+e\)
is greater than the bottom degree \(e\). Its functional on
\(M_l(e)\) comes, after the grading shift, from the functional of
\(z_{e-1}\) on \(M_l(e-1)\). Vanishing of \(z_{e+1}\) says that
it annihilates \(d_0^{e+1}D_l\), by
\eqref{eq:module-suspension}. It must therefore pair nontrivially
with some floor monomial
\[
d_0^e\prod_{j=1}^{l-1}d_j^{b_j}.
\]
In \(M_l(e-1)\) this monomial is above the floor. Since the lifted
functional is \(A\)-annihilated, Lemma~\ref{lem:milnor-coinvariants}
gives
\[
e\equiv l\pmod2,\qquad b_j\text{ odd for every }j\ge1.
\]
The degree formula and \eqref{eq:last-square} yield
\begin{equation}\label{eq:root-parity}
m-1=2^{l-1}e+
\sum_{j=1}^{l-1}(2^{l-1}-2^{j-1})b_j.
\end{equation}
For \(l\ge2\), only the term with \(j=1\) is odd. For \(l=1\),
the right side is \(e\), which is odd by the congruence. Thus
\(m-1\) is odd.

If \(e=1\), \eqref{eq:last-square} shows that \(d\) is odd. Choose
the leading weight \(2^l\) of \(z_0\). Lemma~\ref{lem:zero-weights}
says that it survives to \(z_1\), with a functional lifted from the
annihilated functional on \(M_l(0)\). Since \(z_2=0\), the same
floor argument applies to this component of \(z_1\), and again
\eqref{eq:root-parity} proves that \(m-1\) is odd.

An odd-degree class cannot be a square. Therefore the suspension-kernel
and primitive-root lemmas imply \(\sigma_*w\ne0\) whenever \(e\ge1\).
\end{proof}

Set \(n=e+1\). The map
\[
f_n:S^{2m-1}\longrightarrow T=QS^n
\]
has zero Hurewicz image, and its previous adjoint has image \(w^2\).
The loop product under \(QS^{n-1}\simeq\Omega QS^n\) agrees with
the additive Pontryagin product. Thus Section~\ref{sec:cones} applies
with \(x=w\). The nonzero functional
\begin{equation}\label{eq:last-functional}
\lambda=\langle -,\sigma_*w\rangle
\quad\text{on }\bigl(Q\widetilde H^*(QS^n)\bigr)^m
\end{equation}
gives \(\Sq^m\otimes v\ne0\) for every monomial \(v\) on which
it has value one.

We first deduce that \(m\) is even. If \(m\) were odd, then
\(m\ge3\). Lemma~\ref{lem:algebra-obstruction} excludes a nonzero
pairing on every \(M_k(n)\) with \(k\ge1\). A bottom pairing
requires \(m=n\); on \(M_0(n)=\Sigma^n\F_2\), the balanced tensor
product is the shift of \(I/I^2\). But
\(\Sq^m=\Sq^1\Sq^{m-1}\in I^2\), contradicting
Corollary~\ref{cor:tensor-detection}. Since \(\lambda\ne0\), this
rules out odd \(m\).

\begin{proposition}[The positive terminal stage]\label{prop:hopf-case}
If \(e\ge1\), then \(\alpha\) has nonzero mod-2 stable Hopf
invariant and \(d\in\{1,3,7\}\).
\end{proposition}

\begin{proof}
Suppose that \(\sigma_*w\) has a nonzero component of weight
\(2^k\) with \(k\ge1\). Its degree-\(m\) functional on
\(M_k(n)\) is lifted from the \(A\)-annihilated functional of
\(w\) on \(M_k(e)\). The nonzero tensor supplied by
\eqref{eq:last-functional} therefore satisfies all the hypotheses of
Lemma~\ref{lem:algebra-obstruction}. That lemma forces \(n=1\),
contrary to \(n=e+1\ge2\). Hence \(\lambda\) is nonzero on the
bottom summand, giving
\[
m=n=e+1,\qquad d=e.
\]
Lemma~\ref{lem:cone-square} now says that the extension of the
weight-one bottom class in \(C_{f_n}\) has nonzero \(\Sq^n\), equal
to the top relative class.

Let \(\epsilon:\Sigma^\infty QS^n\to\Sigma^nS\) be the stable
counit and define
\[
\alpha_n=\epsilon\circ\Sigma^\infty f_n:
\Sigma^{2n-1}S\longrightarrow\Sigma^nS.
\]
This is the stable adjoint representing \(\alpha\) in stem
\(d=n-1\). The counit is the weight-one projection, by
Proposition~\ref{prop:free-ring}. The commutative square with identity
on the attaching sphere induces a map
\[
\Sigma^\infty C_{f_n}\longrightarrow C_{\alpha_n}
\]
to the stable two-cell cofiber. Its pullback carries the bottom class
to the weight-one bottom class and the top relative class to the top
relative class. Both top relative classes inject in mod-2 cohomology.
Naturality therefore makes \(\Sq^n\) nonzero in \(C_{\alpha_n}\).
This is nonzero mod-2 stable Hopf detection of \(\alpha\).

Adams's Hopf invariant one theorem, in its stable two-cell form
\cite[Theorem~1.1.1(d)]{Adams1960}, permits \(n=1,2,4,8\).
Since \(d=n-1>0\), we obtain \(d=1,3,7\).
\end{proof}

\begin{proposition}[The terminal stage zero]\label{prop:kervaire-case}
If \(e=0\), then for some \(j\ge1\),
\[
d=2^{j+1}-2,
\]
and \(\alpha\) has Kervaire invariant one.
\end{proposition}

\begin{proof}
\emph{Leading weight.}
Here \(n=1\), and we have proved that \(m\) is even. Thus
\(|w|=m-1\) is odd. Choose the leading weight \(2^k\) of \(w\)
in the level-zero expansion. It has \(k\ge1\), because \(w\) has
positive degree. By Lemma~\ref{lem:zero-weights}, this weight survives
in \(\sigma_*w\), with functional lifted from an annihilated
functional on \(D_k=M_k(0)\). Combining
Corollary~\ref{cor:tensor-detection} with
Lemma~\ref{lem:algebra-obstruction} gives
\[
k=1,\qquad m=2^j\quad\text{for some }j\ge1.
\]
Consequently \(z_0=w^2\) has leading weight four: its leading term
is the square of the nonzero weight-two term, and Frobenius is
injective in the associated free polynomial homology algebra. Also
\(d=2m-2=2^{j+1}-2\).

\emph{Adams filtration.}
We bound the Adams filtration of the given element.
Let \(\mathcal I\) be the augmentation ideal of
\(\Sigma^\infty_+QS^0\). Kuhn's lifting theorem
\cite[Theorems~1.2--1.3]{Kuhn2018}, applied to
\(X=S\) and the coefficient ring \(H\F_2\), says that an element
of mod-2 Adams filtration at least \(s_0\) has Hurewicz image lifted
from \(H\F_2\wedge\mathcal I^{2^{s_0}}\), where
\(\mathcal I^{2^{s_0}}\) is the corresponding derived augmentation-ideal power of
\(\Sigma^\infty_+QS^0\). Its hypotheses hold: both spectra are
connective and \(\pi_0S\to\pi_0H\F_2\) is onto. At the prime two
the localization away from \((p-1)!\) in that theorem is vacuous.
The theorem does not require \(QS^0\) to be connected.

To use this with the leading weight, take the augmented algebra map
\[
\Sigma^\infty_+QS^0\longrightarrow\mathbb P_{\le7}(S\{u\})
\]
from \eqref{eq:weight-projections}, and let \(J\) denote the
positive ideal in its target. Regard \(J\) as a nonunital algebra
in the full category of nonnegatively graded spectra, with ordinary
Day convolution and with \(J_r=0\) for \(r=0\) and \(r\ge8\).
Using the reduced commutative operad \(\mathrm{Com}\), whose arities
are positive, its derived eighth power is the bar construction
\[
J^8=B(\mathrm{Com}_{\ge8},\mathrm{Com},J);
\]
see \cite[Definition~2.7]{Kuhn2018} and
\cite[Definition~3.1 and Example~3.2]{KuhnPereira2017}. Compute this
derived bar construction with a free resolution in
strictly positive weights. Every simplicial summand then has at least
eight inputs of weight at least one; no nullary operation can remove
such an input.
Its total weight is therefore at least eight.
The canonical map \(J^8\to J\) preserves weight and is zero, since
\(J\) has only weights one through seven. The direct-sum functor
from the full graded category is strong symmetric monoidal and
preserves realizations and homotopy orbits. It therefore carries the
graded free resolution and bar construction to the corresponding
ungraded derived construction, with the same zero canonical map.
Extension of coefficients to \(H\F_2\) also preserves this
construction \cite[Theorem~3.4(d)]{KuhnPereira2017}; in particular,
the canonical map remains zero after smashing with \(H\F_2\).

By naturality, the Hurewicz image of an element of Adams filtration
at least three would vanish under every positive-weight projection
below eight. The weight-four
image just found is nonzero. Thus \(\alpha\) has Adams filtration
at most two. This uses only a lower bound on possible weights; it does
not identify Adams filtration with an exact weight.

\emph{Kervaire detection.}
Positive stable stems are finite by Serre's finiteness theorem
\cite[Chapter~V, \S3, Proposition~3]{Serre1951}: choose an odd
sphere dimension \(N>d+1\), so that \(\pi_{N+d}(S^N)\) is both
stable and finite. Their odd-primary parts have zero
mod-2 Hurewicz image and zero mod-2 detecting invariants, so we may
discard them and use the convergent mod-2 Adams spectral sequence for
the remaining two-primary group. More precisely, Adams's convergence
theorem identifies its successive filtration quotients with
\(E_\infty\) and its filtration intersection with the prime-to-two
torsion, which has been removed
\cite[Theorem~2.1(i),(v),(vi)]{Adams1958}.
Write \(F^s\pi_d^S\) for the subgroup of Adams filtration at least
\(s\) in this two-primary group. Filtration zero has no positive-stem
sphere class. Filtration one has internal degrees \(2^a\), hence
stems \(2^a-1\); none is the positive even stem in question. In
filtration two, the classical calculation
\cite[Theorem~2.5.1(1)--(2)]{Adams1960} gives basis elements \(h_bh_a\),
where \(b\ge a\ge0\) and \(b\ne a+1\). The identity
\(2^a+2^b=2^{j+1}\) forces \(a=b=j\). Therefore
\[
\Ext_A^{2,2^{j+1}}(\F_2,\F_2)=\F_2\{h_j^2\},
\]
and the given element is detected by the surviving class \(h_j^2\).

Browder's detection theorem identifies the framed Kervaire invariant
of an individual stable sphere element with its detection by
\(h_j^2\). The original framed Arf identification and secondary
pairing occur in \cite[Theorem~3.2, Proposition~7.6 and
Theorems~7.7--7.8]{Browder1969}; the precise elementwise statement is
also recorded in \cite[Theorem~1.1]{BHHM2020}.
This statement includes the low framed cases \(j=1,2,3\).
It applies to the particular \(\alpha\) above, whose nonzero class
in \(F^2\pi_d^S/F^3\pi_d^S\) is represented by the surviving
\(h_j^2\), and gives Kervaire invariant one.

In particular the detector is unchanged by adding an element of
\(F^3\pi_d^S\). Such an element cannot be detected by a
filtration-two class, so the same elementwise theorem gives it
Kervaire invariant zero; additivity of the framed invariant gives
the assertion. This use of Browder's theorem concerns the invariant
of each element, in addition to the existence of a permanent cycle.
\end{proof}

\begin{proof}[Proof of Theorem~\ref{thm:curtis}]
Every positive-degree stable class with nonzero Hurewicz image has a
last nonzero suspension. Propositions~\ref{prop:hopf-case}
and~\ref{prop:kervaire-case} show that it has the stated nonzero
Hopf or Kervaire invariant. These two degree cases are disjoint.

In a Hopf degree, choose the corresponding class among
\(\eta,\nu,\sigma\). In a Kervaire degree, choose any
Kervaire-invariant-one representative \(\theta_j\); its existence
in this case is already supplied by the class under discussion. Each
detecting invariant is additive with values in \(\F_2\).
For the stable Hopf invariant, attach two top cells to the same bottom
sphere by representatives \(\alpha\) and \(\beta\). If \(u\)
is the bottom class and \(a_\alpha,a_\beta\) are the relative top
classes, naturality on the two one-cell cofibers gives
\(\Sq^{d+1}u=H(\alpha)a_\alpha+H(\beta)a_\beta\).
Pinching the attaching sphere induces a map from the cofiber of
\(\alpha+\beta\) to this two-attachment cofiber; it pulls both
relative classes to the single top class. Naturality therefore gives
\(H(\alpha+\beta)=H(\alpha)+H(\beta)\).
For the Kervaire invariant, disjoint union of framed manifolds takes
their quadratic forms to orthogonal sums and adds their Arf
invariants. Subtracting
the chosen representative therefore gives an invariant-zero class.
If this difference had nonzero Hurewicz image, the proved necessity
would force its invariant to be one, a contradiction. Thus its
Hurewicz image is zero.

Every nonzero image is consequently the image of the chosen Hopf or
Kervaire representative in that degree. These representative images
themselves belong to the image, so they span it. No existence assertion
for an additional \(\theta_j\) is used or obtained from this last
step. This proves Curtis's conjecture in all positive degrees.
\end{proof}
